\chapter{Einleitung}
Die fortschreitende Digitalisierung verändert grundlegende Prozesse der Wissensvermittlung und eröffnet neue Möglichkeiten für individualisierte und adaptive Lernumgebungen. Insbesondere der Einsatz Künstlicher Intelligenz (KI) im Bildungskontext gewinnt zunehmend an Bedeutung, da intelligente Systeme Lernende gezielt unterstützen, Lernprozesse personalisieren und den Zugang zu komplexen Inhalten erleichtern können. Vor diesem Hintergrund rücken dialogische Assistenzsysteme auf Basis grosser Sprachmodelle (Large Language Models, LLMs) in den Fokus der Bildungsforschung und -praxis. Durch ihre Faehigkeit, natürliche Sprache zu verstehen und kontextbezogen zu generieren, ermöglichen sie eine interaktive Form der Lernunterstützung, die über klassische, statische Lernressourcen hinausgeht.
% Quellen für den vorigen Absatz?

LLM-basierte Systeme bieten das Potenzial, Lernende bei Verständnisfragen zu unterstützen, Inhalte auf unterschiedlichen Komplexitätsniveaus zu erklären und Wissen kontextsensitiv zu vermitteln \cite{Neumann2024LLMHE,OatesJohnson2024}. Gleichzeitig weisen rein generative Modelle eine zentrale Limitation auf: Ohne Zugriff auf externe Wissensquellen basieren ihre Antworten ausschliesslich auf Trainingsdaten, wodurch ungenaue oder nicht verifizierbare Aussagen entstehen können. Diese sogenannten Halluzinationen stellen insbesondere im Bildungskontext ein erhebliches Problem dar, da die Verlässlichkeit und Nachvollziehbarkeit von Informationen für effektives Lernen essenziell sind \cite{Niu2024RAGTruth,Yu2024RAGEvalSurvey}.

Ein etablierter Ansatz zur Überwindung dieser Limitation ist Retrieval-Augmented Generation (RAG), bei dem ein Sprachmodell mit einem Retrieval-System kombiniert wird, das relevante Informationen aus einer externen Wissensbasis abruft und in die Antwortgenerierung integriert \cite{Lewis2020RAG}. Durch diese Kombination können Antworten faktenbasierter, aktueller und kontextuell präziser gestaltet werden. RAG-basierte Systeme gelten daher als besonders geeignet für den Einsatz in wissensintensiven Domänen wie der Bildung, da sie generative Sprachfähigkeiten mit strukturiertem, domänenspezifischem Wissen verbinden \cite{Lewis2020RAG,SwachaGracel2025}.

% Der Übergang ist irgendwie nicht smooth, eventuell einen Absatz über Einsatz von RAG im Bildungswesen und dann Übergang zu KI-Campus als einen Player im Bildungskontext?
Der Anwendungskontext dieser Arbeit ist der KI-Campus, eine bundesweit zugängliche Lernplattform für Künstliche Intelligenz. Die Plattform stellt frei verfügbare Online-Kurse, Lernmaterialien, Videos und Podcasts für unterschiedliche Zielgruppen bereit, darunter Studierende, Lehrende, Fachkräfte und die allgemeine Oeffentlichkeit. Innerhalb dieser Lernumgebung übernimmt ein RAG-gestützter Chatbot eine zentrale Rolle als dialogischer Assistent. Er unterstützt Nutzende bei der Navigation durch das Kursangebot, beantwortet fachliche Fragen zu Lerninhalten und dient als interaktive Lernhilfe. Damit bildet der Chatbot eine zentrale Schnittstelle zwischen Nutzenden und der umfangreichen Wissensbasis der Plattform.

Obwohl der Chatbot produktiv eingesetzt wird, zeigen Analysen mehrere Optimierungspotenziale. Einschränkungen betreffen insbesondere Antwortqualität, Kontextsensitivität, Robustheit gegenueber unterschiedlichen Eingaben sowie die Konsistenz der Datenverarbeitung. Diese Faktoren beeinträchtigen sowohl die Zuverlässigkeit als auch die Effektivität des Systems als lernunterstuetzendes Werkzeug. Gleichzeitig existiert in der Forschung bislang kein allgemein akzeptierter Standard zur Evaluation von RAG-basierten Chatbots \cite{Yu2024RAGEvalSurvey,SwachaGracel2025}. Die Bewertung solcher Systeme erfordert eine multidimensionale Betrachtung, die sowohl technische Leistungsmerkmale als auch nutzerzentrierte Qualitätsdimensionen einbezieht \cite{Es2023RAGAS,SaadFalcon2024ARES}.

Vor diesem Hintergrund verfolgt das vorliegende Projekt das Ziel, einen produktiv eingesetzten RAG-basierten Chatbot im Bildungskontext systematisch zu evaluieren und gezielt weiterzuentwickeln. Dabei wird ein mehrdimensionaler Evaluationsansatz verfolgt, der automatisierte Qualitätsmetriken mit qualitativen und quantitativen Nutzerevaluationen kombiniert.

Die Arbeit ist entlang der projektlogischen Abfolge von Analyse, Weiterentwicklung und Evaluation strukturiert. Kapitel~\ref{sec:kic} stellt die bestehende Systemarchitektur des KI-Campus-Chatbots dar und analysiert deren funktionale und strukturelle Eigenschaften als Ausgangspunkt der Untersuchung. Kapitel~\ref{sec:projektdesign} beschreibt das Projekt- und Studiendesign, einschliesslich der Zielsetzung, der Forschungsfrage sowie der entwickelten Evaluationsmethodik. Aufbauend darauf widmen sich die folgenden Kapitel der technischen Weiterentwicklung der RAG-Architektur und der Preprocessing-Pipeline, der empirischen Evaluation der Systemversionen sowie der Diskussion der Ergebnisse und daraus abgeleiteten Implikationen für Forschung und Praxis.
% Die "großen" Kapitel würde ich hier noch bennene (Kapitel 4 ..., Kapitel 5 ... usw. aber halt die Unterkapitel nicht)
